\textsc{The anti\hyp{}\YetterDrinfeld{} modules}
of a finite\hyp{}dimensional Hopf algebra
are a module category
over the \YetterDrinfeld{} modules.
Subsequently, they are implemented by a comodule algebra over the Drinfeld double,
see~\cite{Hajac2004a}.
As explained in \cref{sec: pivotality_of_the_df_centre},
we find ourselves in a similar situation.
The anti\hyp{}Drinfeld centre,
our replacement of the anti\hyp{}\YetterDrinfeld{} modules,
is a module category over the Drinfeld centre.

Replacing finite\hyp{}dimensional vector spaces by
a rigid, possibly pivotal, category \(\cat{C}\),
and the underlying Hopf algebra with a Hopf monad \(H\) on \(\cat{C}\),
this section serves to study a Hopf monad \(D(H) \from \cat{C} \to \cat{C}\)
and over it a comodule monad \(Q(H)\from \cat{C} \to \cat{C}\),
which realise the centre and its twisted cousin as their respective modules.
\Bruguieres{} and Virelizier gave a transparent description of \(D(H)\) in~\cite{Bruguieres2012}
by extending results of Day and Street,~\cite{Day2007}.
If \(H\) is the identity functor,
one defines the \emph{central Hopf monad} \(\mathfrak{D}(\cat{C}^H)\) on \(\cat{C}^H\),
with \(\cent*{\cat{C}^H}\) as its \EilenbergMoore{} category.
As an application of Beck's theory of distributive laws,
one obtains the Drinfeld double \(D(H) \from \cat{C} \to \cat{C}\).
We apply the same techniques to define the anti\hyp{}double \(Q(H)\) of \(H\),
whose modules are isomorphic to the ``dual'' of the anti\hyp{}Drinfeld centre \(\QCat*{\cat{C}^H}\).
This approach is best summarised by \cref{fig:twisted_centres_and_monads}.
We obtain a monadic version of \cref{thm: equiv_in_Hopf_setting}.
\begin{reptheorem}{thm: equivalence_in_the_monad_setting}
  Let \(\cat{C}\) be a rigid monoidal category,
  and suppose that \(H \from \cat{C} \to \cat{C}\) is a Hopf monad
  that admits a double \(D(H)\) and anti\hyp{}double \(Q(H)\).
  The following statements are equivalent:
  \begin{thmlist}[noitemsep]
    \item the monoidal unit \(1 \in \cat{C}\) lifts to a module over \(Q(H)\),
    \item there is an isomorphism of comodule monads \(D(H) \cong Q(H)\), and
    \item there is an isomorphism of monads \(Q(H) \cong D(H)\).
  \end{thmlist}
  If \(\cat{C}\) is pivotal with pivotal structure \(\phi\),
  any of the above statements hold if and only if \(H\) and \(\phi\) admit a pair in involution.
\end{reptheorem}
\begin{figure}[htbp]
  \centering\vspace{-0.2cm}% XXX
  \begin{mytikzcd}[column sep={1cm,between origins}, row sep={0.875cm,between origins}, ampersand replacement=\&]
    \& \&
    \cent*{\cat{C}^H}
    \arrow[rrddd, "U^{(Z)}" description, bend left]
    \arrow[lldddd, leftrightarrow, "K^{\mathfrak{D}}" description, bend left=15]
    \& \& \& \&
    \QCat*{\cat{C}^H}
    \arrow[llddd, "U^{(Q)}" description, bend left]
    \arrow[llll, "\textnormal{action}" description, no head, dotted]
    \arrow[rrdddd, leftrightarrow, "K^{\mathfrak{Q}}" description, bend right=15]
    \& \& \\
    \& \& \& \& \& \& \& \& \\
    \& \& \& \& \& \& \& \& \\
    \& \& \& \&
    {\cat{C}^H}
    \arrow[rruuu, "F^{(Q)}" description, bend left]
    \arrow[lluuu, "F^{(Z)}" description, bend left]
    \arrow[lllld, "F^{\mathfrak{D}}" description, bend left=20]
    \arrow[rrrrd, "F^{\mathfrak{Q}}" description, bend left=20]
    \arrow[ddd, "U^H" description, bend left]
    \& \& \& \& \\
    {{\cat{C}^H}}^{\mathfrak{D}}
    \arrow[rrrru, "U^{\mathfrak{D}}" description, bend left=20]
    \arrow[rrdddd, leftrightarrow, "\Sigma^{(D(H))}" description, bend left=15]
    \& \& \& \& \& \& \& \&
    {{\cat{C}^H}}^{\mathfrak{Q}}
    \arrow[llllu, "U^{\mathfrak{Q}}" description, bend left=20]
    \arrow[lldddd, leftrightarrow, "\Sigma^{(Q(H))}" description, bend right=15]
    \\
    \& \& \& \& \& \& \& \& \\
    \& \& \& \&
    \cat{C}
    \arrow[uuu, "F^H" description, bend left]
    \arrow[rrdd, "F^{Q(H)}" description, bend left]
    \arrow[lldd, "F^{D(H)}" description, bend left]
    \& \& \& \& \\
    \& \& \& \& \& \& \& \& \\
    \& \&
    \cat{C}^{D(H)}
    \arrow[rruu, "U^{D(H)}" description, bend left]
    \& \& \& \&
    \cat{C}^{Q(H)}
    \arrow[lluu, "U^{Q(H)}" description, bend left]
    \arrow[llll, "\textnormal{action}" description, no head, dotted]
    \& \&
  \end{mytikzcd}\vspace{-0.1cm}% XXX
  \scaption[-1.5cm]{%
      A cobweb of adjunctions, monads, and various versions of the centre and anti\hyp{}centre.%
      \label{fig:twisted_centres_and_monads}%
  }
\end{figure}

From \cref{thm: equivalence_in_the_monad_setting} we can deduce how pivotal structures on \(\cat{C}^H\) arise from module morphisms
between the central Hopf monad \(\mathfrak{D}\) and the anti\hyp{}central comodule monad \(\mathfrak{Q}\).

\begin{reptheorem}{cor: pivotal_from_central_anti_central}
  Let \(\cat{C}\) be a rigid monoidal category.
  If \(\cat{C}\) admits a central Hopf monad \(\mathfrak{D}(\cat{C})\)
  and an anti\hyp{}central comodule monad \(\mathfrak{Q}(\cat{C})\),
  then it is pivotal if and only if \(\mathfrak{D}(\cat{C}) \cong \mathfrak{Q}(\cat{C})\) as monads.
\end{reptheorem}

\section{Cross products and distributive laws}\label{sec:cross_products_and_distributive_laws}

\textsc{The Hopf monadic description}
of the Drinfeld centre \(\ZCat*{\cat{C}^H}\), for a rigid category \(\cat{C}\) and Hopf monad \(H\) on \(\cat{C}\),
is achieved as a two-step process:
one first finds a suitable monad on \(\cat{C}^H\), which is then ``lifted'' to a monad on \(\cat{C}\).
We shall review this lifting process based on~\cite[Sections~3 and~4]{Bruguieres2012}.

\begin{definition}\label{def: crossed_product_monad}
  Let \(H\from \cat{C} \to \cat{C}\) be a monad
  and \(T \from \cat{C}^H \to \cat{C}^H\) a functor.
  The \emph{cross product} \(T\rtimes H\) of \(T\) by \(H\) is the endofunctor
  \(U^H T F^H \from \cat{C} \to \cat{C}\).
\end{definition}

If \(\stdmonad\) is a monad, then the cross product \(T \rtimes H\) inherits this structure:
the multiplication and unit are given by
\marginnote{%
  Here \(\varepsilon\) and \(\eta\) are the unit and counit of the adjunction \(F^H \adjoint U^H\).%
}
\[
  \mu^{(T \rtimes H)}\from
  U^H T F^H U^H T F^H
  \xrightarrow{U^H T \varepsilon TF^H} U^H T T F^H
  \xrightarrow{U^H \mu^{(T)}F^H} U^H T F^H
\]
and
\[
  \eta^{(T \rtimes H)}\from \Id_{\cat{C}}
  \xrightarrow{\;\;\eta\;\;} U^H F^H
  \xrightarrow{U^H \eta^{(T)}} U^H T F^H.
\]

Given two bimonads \(H \from \cat{C} \to \cat{C}\) and \(B \from \cat{C}^H \to \cat{C}^H\),
iteratively applying
the comultiplication and counit of \(U^H\), \(B\), and \(F^H\)
yields a bimonad structure on \(B \rtimes H \from \cat{C} \to \cat{C}\).
The comultiplication is given by
\[
  U^H_{2; TF^H(\blank), TF^H(\bblank)} \circ U^H T_{2; F^H(\blank), F^H(\bblank)} \circ U^H T F^H_{2; \blank, \bblank},
\]
and for the counit we have
\[
  \varepsilon^{(U^H)} \circ U^H \varepsilon^{(T)} \circ U^H T \varepsilon^{(F^H)}.
\]

Similar considerations imply the following result.

\begin{lemma}\label{lem: cross_product_of_bi/comonads_is_bi/comonad}
  Let \(H \from \cat{C} \to \cat{C}\) and \(B \from \cat{C}^H \to \cat{C}^H\)
  be bimonads which respectively coact on the comodule monads
  \(G \from \cat{M} \to \cat{M}\) and \(C \from \cat{M}^G \to \cat{M}^G\).
  The cross product \(C \rtimes G \from \cat{M} \to \cat{M}\)
  is a comodule monad over \(B \rtimes H\) via the coaction
  \[
    U^G_{\mathsf{a}; C F^G(\blank), B F^H(\bblank)} \circ U^G C_{\mathsf{a}; F^G(\blank), F^H(\bblank)} \circ U^G C F^G_{\mathsf{a}; \blank, \bblank}.
  \]
\end{lemma}

\begin{remark}
  Let \(H \from \cat{C} \to \cat{C}\) and \(B \from \cat{C}^H \to \cat{C}^H\) be monads.
  The question under which conditions
  the modules \(\cat{C}^{B\rtimes H}\) of \(B\rtimes H\) are isomorphic to \({(\cat{C}^H)}^B\)
  is closely related to the theory of distributive laws of \cref{sec:distributive-laws}.
  We may apply the formal theory thereof, see~\cite{Street1972},
  to the bicategory \(\BiCat{OplMon^{opl}}\) of
  monoidal categories,
  oplax monoidal functors,
  and oplax monoidal natural transformations,
  to obtain a description of bimonads and \emph{oplax monoidal distributive laws}%
  —oplax monoidal natural transformations \(\Lambda \from HB \to BH\) between
  bimonads \(H, B \from \cat{C} \to \cat{C}\)
  that are moreover distributive laws, see~\cite{McCrudden2002}.

  Suppose \(\Lambda \from HB \to BH\) to be an oplax monoidal distributive law.
  The comultiplication and counit of the underlying functor
  \(BH \from \cat{C} \to \cat{C}\) turn the composite \(B \circ_{\Lambda} H\) into a bimonad.
  Thus, comodule monads can be intrinsically described in the bicategory that has
  \begin{itemize}[noitemsep]
    \item as objects, pairs \((\cat{M}, \cat{C})\)
    consisting of a right module category \(\cat{M}\) over a monoidal category \(\cat{C}\);
    \item as 1-morphisms, pairs \((G, F)\)
    of a comodule functor \(G\) over an oplax monoidal functor \(F\); and
    \item as 2-morphisms, pairs \((\phi, \psi)\)
    that form a comodule transformation.
  \end{itemize}
\end{remark}

The subsequent results arise immediately from~\cite{Street1972}.

\begin{definition}\label{def: comodule_distributive_law}
  Let \(G, C \from \cat{M} \to \cat{M}\) be two comodule monads
  over the bimonads \(H,B \from \cat{C} \to \cat{C}\), respectively.
  A \emph{comodule distributive law} is a pair of distributive laws
  \(\Omega \from GC \to CG\) and \(\Lambda \from HB \to BH\)
  such that \((\Lambda, \Omega)\) is a comodule natural transformation.
\end{definition}

\begin{proposition}\label{prop: distributive_laws_are_lifts_of_monads}
  Consider two comodule monads \(G, C \from \cat{M} \to \cat{M}\)
  over the bimonads \(H,B \from \cat{C} \to \cat{C}\).
  There exists a bijective correspondence between:
  \begin{thmlist}[noitemsep]
    \item comodule distributive laws \(\big(GC \trightarrow{\;\;\Omega\;\;} CG, HB \trightarrow{\;\;\Lambda\;\;} BH\big)\); and
    \item lifts of \(B\) to a bimonad \(\widetilde{B} \from \cat{C}^H \to \cat{C}^H\)
    together with lifts of \(C\) to a comodule monad \(\widetilde{C} \from \cat{M}^G \to \cat{M}^G\) over \(\widetilde{B}\),
    such that \(B U^H = U^H\widetilde{B}\) as oplax monoidal functors
    and \(C U^G = U^G \widetilde{C}\) as comodule functors.
  \end{thmlist}
\end{proposition}

Let \(\big(GC \trightarrow{\;\;\Omega\;\;} CG, HB \trightarrow{\;\;\Lambda\;\;} BH\big)\) be a comodule distributive law.
The coactions of \(G\) and \(C\) turn \(C \circ_{\Omega} G\) into a comodule monad over \(B \circ_{\Lambda} H\).

\begin{lemma}\label{lem: distributive_laws_and_comparison_functor}
  Suppose \(\Omega \from GC \to CG\) and \(\Lambda \from HB \to BH\) to form a comodule distributive law.
  Then there are equivalences \({(\cat{C}^H)}^{\widetilde{B}^{\Lambda}} \simeq \cat{C}^{B \circ_{\Lambda} H}\) as monoidal categories,
  and \({(\cat{M}^G)}^{\widetilde{C}^{\Omega}} \simeq \cat{M}^{C \circ_{\Omega} G}\) as \(\cat{C}^{B \circ_{\Lambda} H}\)\hyp{}module categories.
\end{lemma}

In fact, by~\cite[Section~4.5]{Bruguieres2012}, the previous results transcend to the Hopf monadic setting.
Let \(B, H \from \cat{C} \to \cat{C}\) be Hopf monads.
If \(\Lambda \from HB \to BH\) is an oplax monoidal distributive law,
then \(B \circ_{\Lambda} H\from \cat{C} \to \cat{C}\)
and the lift \(\widetilde{B}^{\Lambda} \from \cat{C}^H \to \cat{C}^H\) are Hopf monads as well.

\section{Centralisable functors and the central bimonad}\label{section: centralisable_functors_and_the_central_bimonad}

\textsc{The construction of the double} of a Hopf monad \(H\from \cat{C} \to \cat{C}\) given in~\cite{Bruguieres2012} relies on
explicitly describing the left dual of the forgetful functor \(U^{(Z)} \from \ZCat*{\cat{C}^H} \to \cat{C}^H\),
which is realised as a coend.

\begin{definition}\label{def:centraliser}
  Let \(\cat{C}\) be a rigid category and \(T \from \cat{C} \to \cat{C}\) an endofunctor.
  We call \(T\) \emph{centralisable} if the following coend exists for all \(x \in \cat{C}\):
  \[
    Z_T(x) \defeq \int^{y \in \cat{C}} \ld{Ty} \otimes x \otimes y.
  \]
\end{definition}

\begin{remark}\label{rmk:universal-coaction}
  Any centralisable functor \(T \from \cat{C} \to \cat{C}\) admits a \emph{universal coaction}
  \[
    \chi_{x,y} \defeq (\id_{Ty} \otimes \mathrm{copr}_{y,x}) \circ (\coev^l_{Ty} \otimes\, \id_{x \otimes y}),
    \qquad \text{ for } x, y \in \cat{C},
  \]
  which is natural in both variables.
  We call the pair \((Z_T, \chi)\) a \emph{centraliser} of \(T\).
\end{remark}

\begin{remark}\label{rmk:universal-coaction-graphically}
  Graphically, we represent the universal coaction as follows:
  \[
    \tikzfig{universal-coaction}
  \]
  For all \(f\from x\to x'\) and \(g\from y\to y'\), the naturality condition equates to
  \[
    \tikzfig{universal-coaction-naturality}
  \]
\end{remark}

Universal coactions have a certain universal property that will be vital in the rest of this chapter%
—the \emph{extended factorisation property}.
In particular, it provides us with a potent tool for constructing bi- and comodule monads.

\begin{lemma}[{\cite[Lemma~5.4]{Bruguieres2012}}]\label{lem: extended_factorisation_property}
  Let \((Z_T, \chi)\) be the centraliser of a functor \(T \from \cat{C} \to \cat{C}\)
  and suppose that \(L, R \from \cat{D} \to \cat{C}\) are two functors.
  For any \(n \in \nat\), \(y \in \cat{D}\), and \(x_1, \dots, x_n \in \cat{C}\), and any natural transformation
  \[
    \phi_{y, x_1, \dots, x_n}\from Ly \otimes  x_1 \otimes \dots \otimes x_n \to Tx_1 \otimes \dots \otimes Tx_n \otimes Ry,
  \]
  there exists a unique natural transformation \(\nu \from Z_T^n L \nt R\), satisfying
  \[
    \tikzfig{univ-coaction-univ-prop}
  \]
\end{lemma}

\begin{example}\label{ex:unit-and-multiplication-from-universal-prop-of-universal-coaction}
  Let \(T \from \cat{C} \to \cat{C}\) be an oplax monoidal functor with centraliser \((Z_T, \chi)\).
  For all \(x \in \cat{C}\),
  define the unit of \(Z_T\) by
  \[
    \eta^{(Z_T)}_x \from x \xrightarrow{\ \chi_{x, 1}\ } T1 \otimes Z_T x \xrightarrow{\ T_0 \otimes Z_T x} Z_T x.
  \]

  By \cref{lem: extended_factorisation_property},
  we can derive a unique multiplication \(\mu^{(Z_T)} \from Z_T^2 \nt Z_T\) from the comultiplication of \(T\).
  \[
    \tikzfig{mult-by-univ-fact}
  \]
\end{example}

The following result is due to Day and Street,~\cite[191--192]{Day2007},
see also~\cite[Theorem~5.6]{Bruguieres2012} for a proof.

\begin{lemma}\label{lem: oplax_centralisable_leads to_monad}
  The centraliser \((Z_T, \chi)\)  of an oplax monoidal endofunctor \(T\) on \(\cat{C}\)
  is a monad with multiplication and unit as given in
  \cref{ex:unit-and-multiplication-from-universal-prop-of-universal-coaction}.
\end{lemma}

In the proof of \cref{lem: oplax_centralisable_leads to_monad} given in \cite[Theorem~5.6]{Bruguieres2012},
the authors further consider \(T \from \cat{C} \to \cat{C}\) to be equipped with a Hopf monad structure and show that in this case \(Z_T\) is a Hopf monad as well.
The extended factorisation property given in \cref{lem: extended_factorisation_property}
reconstructs a comultiplication on \(Z_T\)
from a twofold application of the universal coaction and the multiplication of \(T\):
\begin{equation}\label{eq:comult-of-centraliser}
  \tikzfig{comult-of-centraliser}
\end{equation}

Likewise, the unit of \(T\) induces a counit on \(Z_T\) via
\[
  \tikzfig{counit-of-centraliser}
\]
A direct computation verifies that the centraliser \(Z_T\) is a bimonad as well.
For the construction of left and right antipodes, see~\cite[Theorem~5.6]{Bruguieres2012}.

\begin{remark}\label{rem: twisted_centre_of_Hopf_monad_is_rigid}
  We think of \(\ZCat[H]{C}\) as the centre of an oplax bimodule category as stated
  in \cref{rem: other_variants_of_twisted_centres},
  see also~\cite[Section 5.5]{Bruguieres2007}.
  Objects in \(\ZCat[H]{C}\) are pairs \((x, \sigma_{x, \blank})\),
  where \(x \in \cat{C}\) and \(\sigma_{x,\blank} \from x \otimes \blank \nt H(\blank) \otimes x\)
  is a natural transformation, satisfying for all \(x,y,z \in \cat{C}\)
  \begin{align*}
    (H_{2;y,z} \otimes x) \circ \sigma_{x, y \otimes z} &= (Hy \otimes \sigma_{x,z}) \circ (\sigma_{x,y} \otimes z), \\
    (H_0 \otimes x) \circ \sigma_{x,1} &= \id_{x}.
  \end{align*}
  Analogous to the centres studied before,
  the arrows in \(\ZCat[H]{C}\) are those morphisms of \(\cat{C}\)
  that commute with the half-braidings.
  As shown in~\cite[Proposition~5.9]{Bruguieres2012},
  the structure morphisms of a Hopf monad \(H \from \cat{C} \to \cat{C}\)
  can be used to define a rigid structure on \(\ZCat[H]{C}\).
  For example, the tensor product of
  \((x, \sigma_{x,\blank}),\, (y, \sigma_{y,\blank}) \in \ZCat[H]{C}\) is \(x\otimes y \in \cat{C}\),
  together with the half-braiding
  \[
    \tikzfig{tensor-product-with-hopf-monad}
  \]
  For a treatment of centres twisted by (op)lax monoidal functors,
  see~\cite{femić23:categ-yetter}.
\end{remark}

Since centralisers of Hopf monads are Hopf monads themselves,
their modules implement the twisted centres
of \cref{rem: twisted_centre_of_Hopf_monad_is_rigid} as a rigid category.
This is proven in~\cite[Theorem~5.12 and Corollary~5.14]{Bruguieres2012}.

\begin{proposition}\label{prop: central_bimonad_Drinfeld_centre}
  Suppose \(H \from \cat{C} \to \cat{C}\) to be a centralisable Hopf monad.
  The modules \(\cat{C}^{Z_H}\) of its centraliser \((Z_H, \chi)\)
  are isomorphic as a rigid category to \(\ZCat[H]{C}\).
\end{proposition}

Applying the above proposition to the identity functor \(\Id \from \cat{C} \to \cat{C}\),
we obtain a Hopf monadic description of the Drinfeld centre \(\ZCat{C}\) of a rigid category \(\cat{C}\).
The terminology of our next definition is due to Shimizu, see~\cite{Shimizu2017}.

\begin{definition}\label{def: central_bimonad}
  Let \(\cat{C}\) be a monoidal category,
  and let the identity functor on \(\cat{C}\) be centralisable with centraliser \((Z, \chi)\).
  The \emph{central Hopf monad of \(\cat{C}\)}
  is
  \(\mathfrak{D}(\cat{C}) \defeq Z_{\mathrm{Id}} \from \cat{C} \to \cat{C}\).
  We denote its Eilenberg–Moore by \(\cat{C}^{\mathfrak{D}}\).
\end{definition}

An important step in proving \cref{prop: central_bimonad_Drinfeld_centre}
is determining an inverse to the comparison functor \(K^{Z_T}\from \ZCat[T]{C} \to \cat{C}^{Z_T}\).
This construction will also play a substantial role in our monadic description of the anti\hyp{}Drinfeld centre,
so we recall it in its full generality.
Let \(T \from \cat{C} \to \cat{C}\) be a centralisable oplax monoidal endofunctor
with \((Z_T, \chi)\) as its centraliser.
To every \(Z_T\)\hyp{}module \((m, \nabla_m)\),
we associate a half-braiding \(\sigma_{m, \blank} \from m \otimes \blank \nt T(\blank) \otimes m\).
For any \(x \in \cat{C}\), it is given by the composition
\begin{equation}\label{eq: free_half_braiding}
  \sigma_{m, x} \colon m \otimes x \xrightarrow{\ \chi_{m,x}\ } Tx \otimes Z_T x \xrightarrow{Tx \otimes \nabla_m} Tx \otimes m.
\end{equation}
This yields a functor \(E^{Z_T}\from \cat{C}^{Z_T} \to \ZCat[T]{C}\),
which is the identity on morphisms and on objects is given by
\begin{equation}\label{eq: inverse_comparison_functor}
  E^{Z_T}(m, \nabla_{m}) = (m, \sigma_{m,\blank}), \qquad \text{ for all }(m, \nabla_{M}) \in \cat{C}^{Z_T}.
\end{equation}

Conversely, to every object \((m, \sigma_{m, \blank}) \in \ZCat[T]{C}\)
we may assign a \(Z_T\)\hyp{}module,
whose action \(\nabla_m\) is,
due to \cref{lem: extended_factorisation_property},
uniquely defined by
\[
  \tikzfig{action-via-univ-coaction}
\]
This yields the comparison functor \(K^{Z_T}\from \ZCat[T]{C} \to \cat{C}^{Z_T}\).

\begin{remark}\label{rem: adjoint_from_inverse_of_the_comparison_functor}
  Suppose \(T \from \cat{C} \to \cat{C}\) to be a centralisable oplax monoidal endofunctor
  with  \((Z_T, \chi)\) as its centraliser.
  Denote the free functor of the \EilenbergMoore{} adjunction of \(Z_T\) by \(F^{Z_T} \from \cat{C} \to \cat{C}^{Z_T}\).
  The composition
  \[
    \cat{C} \trightarrow{\,F^{Z_T}\,} \cat{C}^{Z_T} \trightarrow{\,E^{Z_T}\,}  \ZCat[T]{C}
  \]
  defines a left adjoint of the forgetful functor \(U^{(T)} \from \ZCat[T]{C} \to \cat{C}\).
\end{remark}

In fact, the central adjunction \(F^{(T)} \adjoint U^{(T)}\) is monadic.

\begin{proposition}[{\cite[Theorem~5.12]{Bruguieres2012}}]\label{prop: inverse_comparison_functor}
  Let \((Z_T, \chi)\) be a centraliser of the oplax monoidal endofunctor
  \(T \from \cat{C} \to \cat{C}\).
  The comparison functor \(K^{Z_T} \from \ZCat[T]{C} \to \cat{C}^{Z_T}\)
  is an isomorphism of categories with inverse \(E^{Z_T} \from \cat{C}^{Z_T} \to \ZCat[T]{C}\).
\end{proposition}

\section{Centralisers and comodule monads}

\textsc{We will now apply the methods of} \Bruguieres{} and Virelizier to twisted centres,
for the purpose of obtaining a comodule monad that implements the anti\hyp{}Drinfeld centre.
Hereto, we need a generalised version of the concept of modules over a monad.
Our approach is based on~\cite{Mesablishvili2011}.

\begin{definition}\label{def: oplax_right_action}
  Let \((B, \mu, \eta) \from \cat{C} \to \cat{C}\) be a bimonad
  and \(F \from \cat{C} \to \cat{D}\) an oplax monoidal functor.
  An \emph{oplax monoidal right action} of \(B\) on \(F\) is
  an oplax natural transformation \(\alpha\from FB \nt F\),
  such that the following diagrams commute:
  \[
    \begin{mytikzcd}[ampersand replacement=\&]
      {FBB} \& FB \& F \& FB \\
      FB \& F \&\& F
      \arrow["F\mu", from=1-1, to=1-2]
      \arrow["\alpha", from=1-2, to=2-2]
      \arrow["{\alpha B}"', from=1-1, to=2-1]
      \arrow["\alpha"', from=2-1, to=2-2]
      \arrow["F\eta", from=1-3, to=1-4]
      \arrow["{\id_F}"', from=1-3, to=2-4]
      \arrow["\alpha", from=1-4, to=2-4]
    \end{mytikzcd}
  \]

  Similarly, one defines oplax monoidal left actions.
\end{definition}

\begin{remark}\label{rmk:oplax-nat-trafo-as-monoid-action}
  Recall from \cref{rmk:bimonad-in-bicategorical-language} that a bimonad \(B\) is a monoid in the monoidal category \(\BiCat{OplMon^{opl}}(\cat{C},\cat{C})\) of
  oplax monoidal endofunctors on \(\cat{C}\),
  with oplax monoidal natural transformations between them.
  There is an obvious right action of this category on
  \(\BiCat{OplMon^{opl}}(\cat{C},\cat{D})\)
  given by composition of functors.
  In this context,
  an oplax monoidal right action of a bimonad \(B\from \cat{C}\to \cat{C}\)
  is the same as an
  \(\BiCat{OplMon^{opl}}(\cat{C},\cat{D})\)\hyp{}module of \(B\).
\end{remark}

A prime example of an oplax monoidal left action is given
by the forgetful functor \(U^B \from \cat{C}^B \to \cat{C}\)
of a bimonad \(B \from \cat{C} \to \cat{C}\) together with the action \(\nabla = U^B \varepsilon \from B U^B \to U^B\),
see \cref{rmk:diagrammatic-algebras}.

\begin{hypothesis}
  To keep our notation concise, in the following we
  fix an oplax monoidal functor \(L \from \cat{C} \to \cat{C}\)
  with an oplax right action \(\alpha \from LB \nt L\) by a bimonad \(B\from \cat{C} \to \cat{C}\)
  and assume that \(L\) and \(B\) are centralisable.
  Their centralisers will be denoted by \((Z_L, \xi)\) and \((Z_B, \chi)\), respectively.
\end{hypothesis}

We think of \(\ZCat[B]{C}\) as a more general version of the Drinfeld centre
which is supposed to act on \(\ZCat[L]{C}\) from the right.
To emphasise this, and in line with the colouring scheme of \cref{sec: pivotality_of_the_df_centre}, we use black for objects in \(\cat{C}\) or its generalised Drinfeld centre and blue for objects in \(\ZCat[L]{C}\).

Consider objects \((m, \sigma_{m,\blank}) \in \ZCat[L]{C}\) and \((x, \sigma_{x,\blank}) \in \ZCat[B]{C}\).
The action of \(B\) on \(L\) and the half-braidings of \(m\) and \(x\) yield
a natural transformation
\[
  \sigma_{m \otimes x, y} \from m \otimes x \otimes y \to Ly \otimes m \otimes x,
\]
which we can depict graphically by
\begin{equation}\label{eq:right-action-of-gen-centre}
  \tikzfig{right-action-of-gen-centre}
\end{equation}

\begin{lemma}\label{lem: right_action_through_centralisers}
  The centre \(\ZCat[B]{C}\) acts on \(\ZCat[L]{C}\) from the right
  by tensoring the underlying objects
  and gluing together the half-braidings
  as in \cref{eq:right-action-of-gen-centre}.
  With respect to this action,
  the forgetful functor \(U^{(L)} \from \ZCat[L]{C} \to \cat{C} \)
  is a strict comodule functor over \(U^{(B)} \from  \ZCat[B]{C} \to \cat{C}\).
\end{lemma}
\begin{proof}
  We proceed as in~\cite[Proposition~5.9]{Bruguieres2012}.
  Fix objects \((m, \sigma_{m, \blank} ) \in \ZCat[L]{C}\) and \((x, \sigma_{x,\blank}) \in \ZCat[B]{C}\).
  The compatibility of the half-braiding of \(m \otimes x\)
  with the unit of \(\cat{C}\) is a short computation:
  \[
    \tikzfig{half-braiding-and-unit}
  \]

  Similarly, we verify the hexagon axiom:
  \[
    \tikzfig{half-braiding-and-hexagon}
  \]

  The compatibility of the action \(\alpha \from LB \nt L\) with the multiplication and unit of \(B\)
  asserts that \(\ZCat[L]{C}\) is a right \(\ZCat[B]{C}\)\hyp{}module category.
  By construction,
  for all \((m, \sigma_{m, \blank} ) \in \ZCat[L]{C}\) and \((x, \sigma_{x,\blank}) \in \ZCat[B]{C}\)
  we have
  \[
    U^{(L)}\big( (m, \sigma_{m,\blank}) \ract (x, \sigma_{x,\blank}) \big)
    \,=\, m \otimes x \,=\,
    U^{(L)}(m, \sigma_{m,\blank}) \otimes U^{(B)} (x, \sigma_{x,\blank}).
  \]
  Thus, \(U^{(L)}\) is a strict comodule functor over \(U^{(B)}\).
\end{proof}

\begin{notation}
  We extend our colouring scheme to universal coactions:
  \[
    \tikzfig{colouring-scheme-universal-coactions}
  \]
\end{notation}

\begin{remark}
  The identification of \(\cat{C}^{Z_B}\) and \(\cat{C}^{Z_L}\) with the generalised Drinfeld centre
  and its twisted cousin suggest that \(Z_L\) is a comodule monad over \(Z_B\).
  In analogy with \cref{eq:comult-of-centraliser}
  we define the coaction \(Z_{L;\mathsf{a}}\) of \(Z_L\) by
  \begin{equation}\label{eq:coaction-anti-central-monad}
    \tikzfig{coaction-anti-central-monad}
  \end{equation}
\end{remark}

\begin{proposition}\label{prop: construction_anti_central_comonad}
  Let \(\alpha \from LB \nt L\) be an oplax monoidal right action of a bimonad \(B \from \cat{C} \to \cat{C}\) on an oplax monoidal functor \(L \from \cat{C} \to \cat{C}\).
  Suppose furthermore that the centralisers \((Z_L, \xi)\) of \(L\)  and  \((Z_B, \chi)\) of \(B\) exist.
  The coaction of \cref{eq:coaction-anti-central-monad} turns \(Z_L\) into a comodule monad over \(Z_B\) such that \(\cat{C}^{Z_L}\) is isomorphic as a right module category over \(\cat{C}^{Z_B}\) to \(\ZCat[L]{C}\).
\end{proposition}
\begin{proof}
  By \cref{rem: adjoint_from_inverse_of_the_comparison_functor}
  and \cref{prop: inverse_comparison_functor},
  we have monadic adjunctions
  \[
    \adj{F^{(B)}}{U^{(B)}}{\cat{C}}{\ZCat[B]{C}}
    \qquad \text{ and } \qquad
    \adj{F^{(L)}}{U^{(L)}}{\cat{C}}{\ZCat[L]{C}}
  \]
  that, due to~\cite[Remark~5.13]{Bruguieres2012},
  give rise to the bimonad \(Z_B\) and monad \(Z_L\).
  \Cref{lem: right_action_through_centralisers} shows that
  \(U^{(L)}\) is a strict comodule functor over \(U^{(B)}\).


  By  \cref{thm:comodule-monad-reconstruction},
  see also \cref{rem:comodule-monads-induce-module-categories,ex:comodule-monad-over-bimonad},
  \(Z_L\) is a comodule monad over \(Z_{B}\);
  the coaction \(\lambda \from Z_L(\blank \otimes \bblank) \to Z_L(\blank) \otimes Z_B(\bblank)\) implementing the action of \(\cat{C}^{Z_B}\) on \(\cat{C}^{Z_L}\) is for all \(x, y \in \cat{C}\) given by
  \[
    \lambda_{x,y}\from Z_L(x \otimes y)
    \xrightarrow{Z_L(\eta^{(Z_L)}_{x} \otimes \eta^{(Z_B)}_{y})} Z_L(Z_L x \otimes Z_B y)
    \xrightarrow{\nabla_{Z_L x \otimes Z_B y}} Z_L x \otimes Z_B y,
  \]
  where \(\nabla\) is defined as in \cref{eq:action-as-nat-trafo}.

  By using the relation between universal coactions and half-braidings,
  explained in \cref{eq: free_half_braiding},
  and applying the hexagon identity we compute \cref{fig:lambda-is-really-qa}.
  The uniqueness of universal coactions implies that \(\lambda = Z_{L;\mathsf{a}}\).
  \begin{figure}[htbp]
    \centering
    \vspace{-0.2cm}\tikzfig{lambda-is-really-qa}\vspace{-\onelineskip}% XXX
    \scaption[-.5cm]{%
      The arrows \(\lambda\) and \(Z_{L;\mathsf{a}}\) satisfy the same universal property.%
      \label{fig:lambda-is-really-qa}%
    }
  \end{figure}

  It remains to show that \(\cat{C}^{Z_L}\) and \(\ZCat[L]{C}\) are isomorphic as modules over \(\cat{C}^{Z_B}\).
  By~\cref{prop:comparison-functor-comodule-monad-comodule}, the comparison functor
  \(K^{Z_B} \from \ZCat[B]{C} \to \cat{C}^{Z_B}\) is strong monoidal
  and \(K^{Z_L} \from \ZCat[L]{C} \to \cat{C}^{Z_L}\) is a strong comodule functor over it.
  Furthermore, due to \cref{prop: inverse_comparison_functor},
  both \(K^{Z_B}\) and \(K^{Z_L}\)  admit inverses
  \[
    E^{Z_B} \from \cat{C}^{Z_B} \to \ZCat[B]{C}
    \qquad \text{ and } \qquad
    E^{Z_L} \from \cat{C}^{Z_L} \to \ZCat[L]{C}.
  \]
  Using that \(E^{Z_B}\) is monoidal, we identify the right action of \(\ZCat[B]{C}\) on \(\ZCat[L]{C}\)
  with a right action \(\bract\from\ZCat[L]{C} \times \cat{C}^{Z_B} \to \ZCat[L]{C} \) of \(\cat{C}^{Z_B}\) by setting
  \[
    \ZCat[L]{C} \times \cat{C}^{Z_B} \trightarrow{\id \times E^{Z_B}}
    \ZCat[L]{C} \times \ZCat[B]{C} \trightarrow{\ \ \ract \ \ }
    \ZCat[L]{C}.
  \]
  For any \(m \in \ZCat[L]{C}\) and \(x \in\ZCat[L]{C}\) we have
  \[
    K^{Z_L} (m \bract x)
    =
    K^{Z_L} (m \ract E^{Z_B}x)
    \trightarrow{K^{Z_L}_{\mathsf{a};m, E^{Z_B}x}} K^{Z_L}m \ract K^{Z_B}E^{Z_B}x = K^{Z_L}m \ract x,
  \]
  and hence \(K^{Z_L} \from \ZCat[L]{C} \to \cat{C}^{Z_L}\) is an isomorphism of module categories.
\end{proof}

\begin{example}
  Let \(\cat{C}\) be a rigid monoidal category, and
  let \((Z_L, \xi)\) and  \((Z_B, \chi)\) be
  the centralisers of \(\lbiDual{\!(\blank)}\from \cat{C}\to \cat{C}\) and \(\Id_{\cat{C}}\), respectively.
  Then there exists a trivial right action of \(\Id_{\cat{C}}\) on \(\lbiDual{\!(\blank)}\):
  \[
    \id_x\from \lbiDual{\!(\Id_{\cat{C}} (x))} \to \lbiDual{\!x}, \qquad \qquad \text{for all } x \in \cat{C}.
  \]
  This action turns \(Z_L\) into a comodule monad over \(Z_B\),
  and its modules \(\cat{C}^{Z_L}\) become isomorphic to \(\QCat{C}\) as a \(\cat{C}^{Z_B}\)\hyp{}module category.
\end{example}

Recall that,
by \cref{rem: right_vs_left_twisted_centres},
we can identify \(\QCat{C}\) with \({\ACat*{\cat{C}^{\op,\tensorop}}}^{\op}\).

\begin{definition}\label{def: anti_central_comodule_monad}
  Assume \(\lbiDual{\!(\blank)},\, \Id_{\cat{C}}\from \cat{C} \to \cat{C}\) to admit centralisers \((Z_L, \xi)\) and \((Z_B, \chi)\).
  We call \(\mathfrak{Q}(\cat{C}) \defeq Z_L\) the \emph{anti\hyp{}central comodule monad of \(\cat{C}\)}.
\end{definition}

\section{The Drinfeld and anti-Drinfeld double of a Hopf monad}\label{subsec: Drinfeld_anti-Drinfeld_double}

\textsc{We are now able to untangle}
the relationship between the various different categories and adjunctions in \cref{fig:twisted_centres_and_monads}.

\begin{hypothesis}
  For the rest of this chapter, fix a Hopf monad \(H\) on a rigid category \(\cat{C}\),
  together with an oplax monoidal endofunctor \(L\) on \(\cat{C}^H\),
  a bimonad \(B\) on \(\cat{C}^H\),
  an oplax monoidal right action \(\alpha \from LB \nt B\),
  and assume that the cross products \(B \rtimes H\) and \(L \rtimes H\) have centralisers \((Z_{B \rtimes H}, \nu)\) and \((Z_{L \rtimes H}, \tau)\).
\end{hypothesis}

The following observation%
—which follows by a straightforward calculation—%
extends the action of \(B\) on \(L\) to an action of the respective cross products.

\begin{lemma}\label{lem: cross_products_and_oplax_monoidal_actions}
  The oplax monoidal right action \(\alpha \from LB \nt B\) induces an oplax monoidal action of \(B \rtimes H\) on \(L \rtimes H\) by
  \[
    \tikzfig{oplax-mon-right-action-rtimes}
  \]
\end{lemma}

The next result is a variant of~\cite[Theorem~7.4]{Bruguieres2012}.

\begin{theorem}\label{thm: centralisablity_of_lifts}
  The functors \(B, L\from \cat{C}^H\to \cat{C}^H\) admit centralisers \((Z_B, \chi), (Z_L, \xi)\),
  such that \(Z_B\) lifts \(Z_{B \rtimes H}\) as a bimonad
  and \(Z_L\) lifts \(Z_{L \rtimes H}\) as a comodule monad.
\end{theorem}
\begin{proof}
  By~\cite[Theorem~7.4(a)]{Bruguieres2012},
  there are centralisers \((Z_L, \xi)\) and \((Z_B, \chi)\) of \(L\) and \(B\)
  that, for all \((x, \nabla_x),\, (y, \nabla_y) \in \cat{C}^H\), satisfy
  \begin{align*}
    U^H Z_L (x, \nabla_x) &= Z_{L \rtimes H}x, & U^H\xi_{(x, \nabla_x), (y, \nabla_y)} &= (U^H L\nabla_y \otimes Z_{L \rtimes H}x) \circ \tau_{x,y}, \\
    U^H Z_B (x, \nabla_x) &= Z_{B \rtimes H}x, & U^H\chi_{(x, \nabla_x), (y, \nabla_y)} &= (U^H B\nabla_y \otimes Z_{B \rtimes H}x) \circ \nu_{x,y}.
  \end{align*}
  The second and third part of \emph{ibid} state that \(Z_L\) is a lift of the monad \(Z_{L \rtimes H}\) and \(Z_B\) is a lift of the bimonad \(Z_{B \rtimes H}\).
  It remains for us to show that the coactions of \(Z_L\) and \(Z_{L \rtimes H}\) are compatible with the forgetful functor \(U^H \from \cat{C}^H \to \cat{C}\).
  Fixing objects \((x, \nabla_x),\, (y, \nabla_y) \in \cat{C}^H\) and \(w \in \cat{C}\),
  this follows from \cref{fig:q-and-qh-compatible-with-forgetful}.
  \begin{figure}[htbp]
    \centering
    \vspace{-0.3cm}\mathscale{0.93}{\tikzfig{q-and-qh-compatible-with-forgetful}}% XXX
    \scaption[3.5cm]{%
      The coactions of \(Z_L\) and \(Z_{L \rtimes H}\) are compatible with the forgetful functor.%
      \label{fig:q-and-qh-compatible-with-forgetful}%
    }
  \end{figure}

  The uniqueness property of universal coactions as given in
  \cref{lem: extended_factorisation_property}
  then implies that \(U^H Z_{L;\mathsf{a};(x, \nabla_x), (y, \nabla_y)} = Z_{L \rtimes H;\mathsf{a};x,y}\).
  Since \(U^{H}\from \cat{C}^H \to \cat{C}\) is a strict comodule functor, the claim follows.
\end{proof}

\begin{remark}\label{blackstar}
  The previous theorem together with \cref{lem: cross_product_of_bi/comonads_is_bi/comonad}
  imply that we obtain a comodule monad \(D(L,H) \defeq Z_L \rtimes H\) over \(D(B, H) \defeq Z_B \rtimes H\).
  The correspondence between lifts and monads given in
  \cref{prop: distributive_laws_are_lifts_of_monads} yields a unique comodule distributive law
  \[
    \big(
    H Z_{L \rtimes H}  \trightarrow{\;\Omega\;} Z_{L \rtimes H} H
    ,\;
    H Z_{B \rtimes H} \trightarrow{\;\Lambda\;} Z_{B \rtimes H} H
    \big),
  \]
  such that
  \[
    D(L, H) = Z_{L \rtimes H} \circ_{\Omega} H
    \qquad \text{and} \qquad
    D(B, H) = Z_{B \rtimes H} \circ_{\Lambda} H.
  \]
\end{remark}

\begin{definition}\label{def: double-and-twisted-double}
  We call \(D(B,H)\) and \(D(L,H)\) of \cref{blackstar} the \emph{double} and \emph{twisted double}
  of the pairs \((B,H)\) and \((L,H)\), respectively.
\end{definition}

The relationship between doubles and generalised Drinfeld centres
is studied in~\cite[Proposition~7.5 and Theorem~7.6]{Bruguieres2012}.
Our next result uses the same techniques to prove
how twisted doubles parameterise twisted centres.

\begin{theorem}\label{thm: twisted_doubles_vs_twisted_centres}
  The twisted double \(D(L, H)\) is a comodule monad over \(D(B,H)\),
  and \(\cat{C}^{D(L, H)}\) is isomorphic to \(\ZCat*[L]{\cat{C}^H}\) as a \(\cat{C}^{D(B, H)}\)\hyp{}module category.
\end{theorem}
\begin{proof}
  Since \(Z_L\) is a lift of \(Z_{L \rtimes H}\) as a comodule monad,
  the twisted double \(D(L, H)\) is a comodule monad over \(D(B,H)\).
  By \cref{lem: distributive_laws_and_comparison_functor},
  this implies the existence of an isomorphism
  \(I\from \cat{C}^{D(L, H)} \iso {\left(\cat{C}^H\right)}^{Z_L}\)
  of \(\cat{C}^{D(B, H)}\)\hyp{}module categories.
  Due to the proof of \cref{prop: construction_anti_central_comonad},
  the comparison functor \(K^{Z_L}\from  \ZCat*[L]{\cat{C}^H} \to {\left(\cat{C}^H\right)}^{Z_L}\)
  implements an isomorphism of module categories and the statement follows by considering
  \[
    \cat{C}^{D(L, H)} \trightarrow{\ \ I\ \ }
    {(\cat{C}^H)}^{Z_L} \trightarrow{\,E^{Z_L}\,}
    \ZCat*[L]{\cat{C}^H}.\qedhere
  \]
\end{proof}

The following definition can be understood as an extension of the notion of the anti\hyp{}Drinfeld double given by~\cite{Hajac2004a} to the monadic framework.

\begin{definition}\label{def: double_of_bimonad}
  Let \(H\) be a Hopf monad on a rigid category \(\cat{C}\).
  For \(B \defeq \Id_{\cat{C}^H}\) and \(L \defeq \lbiDual{\!(\blank)} \from \cat{C}^H \to \cat{C}^H\),
  we call \(D(H) \defeq D(B,H)\) and \(Q(H) \defeq D(L,H)\)
  the \emph{Drinfeld} and \emph{anti\hyp{}Drinfeld double} of \(H\), respectively.
\end{definition}

\section{Pairs in involution for Hopf monads}\label{subsec: monadic_perspective_piv_structures}

\textsc{We now consider a Hopf monad}
that admits a double and anti\hyp{}double,
and develop the notion of pairs in involution in this setting.
Classically, these consist of a group-like and character of a Hopf algebra
that implement the square of its antipode by their adjoint actions.

\begin{definition}\label{def: character_and_group-like}
  Let \(H\) be a Hopf monad on a rigid monoidal category \(\cat{C}\).
  A \emph{character} of \(H\) is an \(H\)-algebra \(\beta \defeq (1, \nabla_{\beta}) \in \cat{C}^H\),
  whose underlying object is the monoidal unit \(1 \in \cat{C}\).
\end{definition}

\begin{remark}
  Explicitly, \cref{def: character_and_group-like} says that
  a group-like \(g\) of \(H\) satisfies
  \[
    H_{2;x,y} \circ g_{x\otimes y} = g_x \otimes g_y
    \qquad \text{ and } \qquad
    H_0 \circ g_1= \id_1,\qquad
    \text{for all \(x, y \in \cat{C}\)}.
  \]
\end{remark}

The characters \(\Char(H)\) of a Hopf monad \(H\) on \(\cat{C}\) form a monoid
and, by~\cite[Lemma~3.21]{Bruguieres2007}, the set \(\Gr(H)\) of group-likes bears a group structure.

\begin{example}
  The group-likes of a Hopf monad \(H\) act on it by conjugation.
  We recall this construction based on~\cite[Section~1.4]{Bruguieres2007}.
  Given a natural transformation \(g \from \Id_{\cat{C}} \to H\),
  define the \emph{left} and \emph{right regular action} of \(g\) on \(H\) to be the natural transformations defined by
  \[
    L_{g}\from H \trightarrow{gH} H^2 \trightarrow{\mu^{(H)}} H
    \qquad\text{and}\qquad
    R_{g}\from H \trightarrow{Hg} H^2 \trightarrow{\mu^{(H)}} H.
  \]
\end{example}

\begin{definition}\label{def: adjoint_action}
  Every pair \((g \in \Gr(H),\ \beta \in \Char(H))\)
  of a Hopf monad \(H \from \cat{C} \to \cat{C}\)
  gives rise to natural transformations%
  \marginnote{Recall the definition of \(H_3\) from \cref{rmk:whynotthree}.}
  \begin{align*}
    \Ad_{g} &\defeq L_{g} \circ R_{g^{-1}} \from H \nt H, \\
    \Ad_{\beta} &\defeq (\nabla_{\beta} \otimes H(\blank) \otimes \nabla_{\ld{\!\beta}}) \circ H_{3; 1, \blank ,1} \from H \nt H,
  \end{align*}
  called the \emph{adjoint actions} of \(g\) and \(\beta\) on \(H\), respectively.
\end{definition}

To define pairs in involution, we need an analogue of the square of the antipode of a Hopf algebra.
This notion is developed in~\cite[Section~7.3]{Bruguieres2007}.

\begin{definition} \label{def: s-squared}
  Suppose \(\phi \from \Id_{\cat{C}} \to \lbiDual{\!(\blank)}\) to be a pivotal structure on \(\cat{C}\)
  and let \(H \from \cat{C} \to \cat{C}\) be a Hopf monad.
  The \emph{square of the antipode} of \(H\) is a natural transformation \(S^{2}\from H \nt H\),
  given for all \(x \in \cat{C}\) by
  \[
    S^2_x \defeq \phi^{-1}_{Hx} \circ s^l_{\ld{Hx}} \circ H\ld{s^l_x} \circ H\phi_x,
  \]
  where \(s^l\) is the left antipode of \(H\),
  see \cref{eq: left_and_right_antipode} and~\cite[Section~3.3]{Bruguieres2007}.
\end{definition}

Analogous to the Hopf algebraic case, we state the following:

\begin{definition}\label{def: pair_in_involution_Hopf_monad}
  Let \(H \from \cat{C} \to \cat{C}\) be a Hopf monad,
  and \(\phi \from \Id_{\cat{C}} \to \lbiDual{\!(\blank)}\) a pivotal structure.
  A \emph{pair in involution}
  \((g, \beta) \in \mathrm{PI}_H^{\phi}\)
  of \(H\) and \(\phi\) consists of
  a group-like \(g \in \Gr(H)\) and
  a character \(\beta \in \Char(H)\),
  such that for all \(x \in \cat{C}\)
  \[
    Ad_{g, x} =  Ad_{\beta,x} \circ S^{2}_{x}.
  \]
\end{definition}

To prove that pairs in involution correspond to certain pivotal structures on the Drinfeld centre of \(\cat{C}^H\), we need two technical results.
The first one is classical; for a proof see for example~\cite[Lemmas~1.2 and~1.3]{Bruguieres2007}.

\begin{lemma}\label{lem: sharp-lift}
  Let \(H\) be a monad with canonical forgetful functor \(U^H \from  \cat{C}^H \to \cat{C}\).
  If \(F, G \from \cat{C} \to \cat{D}\) are functors, for some category \(\cat{D}\),
  then there is a canonical bijection
  \[
    {(\blank)}^{\sharp}\from \Nat(F, GH) \to \Nat(FU^H, GU^H), \qquad \alpha \mapsto \alpha^{\sharp},
  \]
  where \(\alpha^{\sharp}_{(m, \nabla_m)} \defeq G\nabla_m \circ \alpha_m\).
\end{lemma}

The next lemma is a variant of~\cite[Lemma~7.5]{Bruguieres2007}.

\begin{lemma}\label{lem: s2-is-conjugate}
  Let \(\phi \from \Id_{\cat{C}} \to \lbiDual{\!(\blank)}\) be a pivotal structure on \(\cat{C}\)
  and \(H \from \cat{C} \to \cat{C}\) a Hopf monad.
  For any group-like \(g \in \Gr(H)\) and character \(\beta \in \Char(H)\) the following are equivalent:
  \begin{thmlist}[noitemsep]
    \item the morphisms \(g\) and \(\beta\) form a pair in involution of \(H\) and \(\phi\); and
    \item the arrow \(\phi g^{\sharp}\in \Nat(U^H, \lld{(\blank)} \circ U^H)\) lifts
    to \(\Nat(\Id_{\cat{C}^H}, \beta \otimes \lbiDual{\!(\blank)}\otimes \ld{\!\beta} )\).
  \end{thmlist}
\end{lemma}
\begin{proof}
  Consider a module \((m, \nabla_m) \in \cat{C}^H\).
  By~\cite[Theorem~3.8(a)]{Bruguieres2007} and the definition of \(S^2\),
  the action on \(\lld{m}\) is given by
  \[
    \nabla_{\lbiDual{m}}
    = \lbiDual{\nabla_m} \circ s^l_{\lbiDual{Hm}} \circ H(\ld{\!{}s^l_m})
    = \phi_m \circ \nabla_m \circ S^2_m \circ H(\phi^{-1}_m),
  \]
  and therefore we have
  \[
    \nabla_{\beta \otimes \lbiDual m \otimes \ld{\!\beta}}
    = (\nabla_{\beta} \otimes \nabla_{\lbiDual{m}} \otimes \nabla_{\ld{\!\beta}}) H_{3; 1,m,1}
    = (\nabla_{\beta} \otimes \phi_m \nabla_m S^2_m H(\phi^{-1}_m) \otimes \nabla_{\ld{\!\beta}}) H_{3; 1,m,1}.
  \]
  By definition, \(\phi g^{\sharp}\) lifts to a natural transformation from \(\Id_{\cat{C}^H}\) to \(\beta \otimes \lbiDual{\!(\blank)} \otimes \ld{\!\beta}\)
  if and only if for any \(H\)\hyp{}module \((m, \nabla_m)\) we have
  \begin{equation} \label{eq: sharpbeta-lifts}
    {(\phi g^{\sharp})}_m \nabla_m = \nabla_{\beta \otimes \lbiDual{m} \otimes \ld{\!\beta}}H \big( {(\phi g^{\sharp})}_m \big).
  \end{equation}

  Let us now successively simplify both sides of this equation.
  Using the naturality of \(g \from \Id_{\cat{C}} \nt H\),
  the fact that \(\nabla_m\) is an action,
  and the definition of \(g^{\sharp}\) as given in \cref{lem: sharp-lift},
  we can rewrite the left hand side as
  \[
    {(\phi g^{\sharp})}_m \nabla_m
    = \phi_m \nabla_m g_m \nabla_m
    = \phi_m \nabla_m H(\nabla_m) g_{Hm}
    = \phi_m \nabla_m \mu^{(H)}_m g_{Hm}.
  \]
  Similarly, we simplify the right-hand side to
  \begin{align*}
    \nabla_{\beta \otimes \lbiDual{m} \otimes \ld{\!\beta}}H\big( {(\phi g^{\sharp})}_m \big)
    & = (\nabla_{\beta} \otimes \phi_m \nabla_m S^2_m H(\phi^{-1}_m) \otimes \nabla_{\ld{\!\beta}}) H_{3; 1,m,1} H\big( {(\phi g^{\sharp})}_m \big) \\
    & = (\nabla_{\beta} \otimes \phi_m \nabla_m  S^2_m H(\phi^{-1}_m) H\big( {(\phi g^{\sharp})}_m\big) \otimes \nabla_{\ld{\!\beta}}) H_{3;1,m,1} \\
    & = (\nabla_{\beta} \otimes \phi_m \nabla_m  S^2_m H(\nabla_m g_m ) \otimes \nabla_{\ld{\!\beta}}) H_{3;1,m,1} \\
    & = (\nabla_{\beta} \otimes \phi_m \nabla_m   H(\nabla_m g_m ) S^2_m \otimes \nabla_{\ld{\!\beta}}) H_{3;1,m,1} \\
    & = \phi_m \nabla_m   H(\nabla_m g_m ) (\nabla_{\beta} \otimes  \id_{Hm}\otimes \nabla_{\ld{\!\beta}}) H_{3;1,m,1}S^2_m \\
    & = \phi_m \nabla_{m}\mu^{(H)}_m H(g_m) \Ad_{\beta,m} S^2_m .
  \end{align*}
  Using the fact that \(\phi\) is an isomorphism,
  \Cref{eq: sharpbeta-lifts} can be restated as
  \begin{align*}
    & & \nabla_m \mu^{(H)}_m g_{Hm} &= \nabla_m\mu^{(H)}_m H(g_m)  \Ad_{\beta,m} S^2_m \\
    \iff & & \nabla_m L_{g, m} &= \nabla_m R_{g, m} \Ad_{\beta, m} S^2_m
  \end{align*}
  By \cref{lem: sharp-lift},
  the above is equivalent to
  \(L_{g, m} = R_{g, m} \Ad_{\beta, m} S^2_m\).
  We conclude the proof by multiplying both sides with \(R_{g^{-1}, m}\).
\end{proof}

\Cref{lem: s2-is-conjugate} leads to an identification of pairs in involution of \(H\) and \(\phi\)
with certain quasi-pivotal structures on \(\cat{C}^H\).

\begin{proposition}\label{prop: piis_are_qpiv}
  Let \(H\) be a Hopf monad on \(\cat{C}\)
  and \(\phi \from \Id_{\cat{C}} \to \lbiDual{\!(\blank)}\) a pivotal structure on \(\cat{C}\).
  Then \((H, \phi)\) admits a pair in involution
  if and only if
  there exists a quasi-pivotal structure on \(\cat{C}^{H}\)
  whose underlying invertible object is a character.
\end{proposition}
\begin{proof}
  We proceed analogous to~\cite[Proposition~7.6]{Bruguieres2007}.
  Suppose \((g, \beta) \in \mathrm{PI}_{H}^{\phi}\).
  By \cref{lem: s2-is-conjugate}, the natural transformation \(\phi g^{\sharp}\) lifts to a natural isomorphism
  \[
    \rho_{\beta,x}\from x \to \beta \otimes \lbiDual x \otimes \ld{\!\beta}, \qquad \text{ for all } x \in \cat{C}^H.
  \]
  Since \(\phi\) is monoidal by definition,
  and \(g^{\sharp}\) is monoidal by virtue of \(g\) being a group-like%
  —see~\cite[Lemma~3.20]{Bruguieres2007}—%
  we obtain a quasi-pivotal structure
  \[
    \rho_{\beta} \from \Id_{\cat{C}^H} \to \beta \otimes \lbiDual{(\blank)}\otimes \ld{\!\beta}.
  \]

  On the other hand, let \((\beta,\rho_{\beta})\) be a quasi-pivotal structure, for \(\beta \in \Char(H)\).
  Since the forgetful functor \(U^H\) is strong monoidal and thus
  \[
    U^H (\beta\otimes \lbiDual{(\blank)}\otimes \ld{\!\beta}) = U^H(\lbiDual{(\blank)}) = \lbiDual{(U^H(\blank))},
  \]
  there exists a monoidal natural transformation
  \[
    \phi^{-1}_{U^H x} \circ U^H(\rho_{\beta,x}) \from U^H x \to U^H x, \qquad \qquad \text{for all } x \in \cat{C}^H.
  \]
  Apply~\cite[Lemma~3.20]{Bruguieres2007} to
  obtain a unique group-like \(g\in \Gr(H)\), with
  \[
    g^{\sharp} = \phi^{-1}_{U^H(x)} \circ U^H(\rho_{\beta,x}).
  \]
  As \(\phi g^{\sharp} = U^H(\rho_{\beta})\) lifts to the quasi-pivotal structure \((\beta,\rho_{\beta})\) on \(\cat{C}^H\),
  \cref{lem: s2-is-conjugate} implies that \((g, \beta) \in \mathrm{PI}_{H}^{\phi}\).
\end{proof}

Let us now study a variant of~\cite[Lemma~2.9]{Bruguieres2007}.

\begin{proposition}\label{prop: functors_as_morphisms_of_monads}
  Let \(C, D \from \cat{M} \to \cat{M}\) be two comodule monads over a bimonad \(B \from \cat{C} \to \cat{C}\).
  There is a bijective correspondence between morphisms of comodule monads \(f \from D \nt C\)
  and strict module functors \(F \from \cat{M}^C \to \cat{M}^D\) with \(U^D F = U^C\).
\end{proposition}
\begin{proof}
  As shown in~\cite[Lemma~1.7]{Bruguieres2007},
  any functor \(F\from \cat{M}^C \to \cat{M}^D\) with \(U^D F = U^C\)
  is induced by a unique  morphism of monads \(f \from D \to C\).
  That is, \(F\) is the identity on morphisms and on objects it is defined by
  \[
    F(m, \nabla_m) = (m, \nabla_m f_m), \qquad \qquad \text{ for all } (m, \nabla_m) \in \cat{m}^C.
  \]

  It remains to show that \(f\) is a morphism of comodules
  if and only if
  \(F\) is a strict \(\cat{C}^B\)\hyp{}module functor.
  Let \((m, \nabla_m)\in \cat{m}^C\) and \((x, \nabla_x) \in \cat{C}^B\).
  We compute
  \begin{align*}
    F\big((m, \nabla_m) \ract (x, \nabla_x)\big)
    &= (m \ract x,\ (\nabla_m \ract \nabla_x) \circ C_{\mathsf{a};m,x} \circ f_{m\ract x}), \\
    F(m, \nabla_m) \ract  (x, \nabla_x)
    &= (m \ract x,\ (\nabla_m \ract \nabla_x) \circ (f_m \ract Bx) \circ D_{\mathsf{a}; m, x}).
  \end{align*}
  According to~\cite[Lemma~1.4]{Bruguieres2007}, these modules coincide if and only if
  \[
    C_{\mathsf{a};m,x} \circ f_{m \ract x} = (f_m \ract Bx) \circ D_{\mathsf{a}; m, x},
  \]
  which is exactly the condition for \(f\) to be a comodule morphism.
\end{proof}

The above result implies the desired monadic version of \cref{thm: equiv_in_Hopf_setting}.

\begin{theorem}\label{thm: equivalence_in_the_monad_setting}
  Let \(\cat{C}\) be a rigid monoidal category,
  and suppose that \(H \from \cat{C} \to \cat{C}\) is a Hopf monad
  that admits a double \(D(H)\) and anti\hyp{}double \(Q(H)\).
  The following statements are equivalent:
  \begin{thmlist}[noitemsep]
    \item\label{big-monadic-equiv-1} the monoidal unit \(1 \in \cat{C}\) lifts to a module over \(Q(H)\);
    \item\label{big-monadic-equiv-2} there is an isomorphism of comodule monads \(D(H) \cong Q(H)\); and
    \item\label{big-monadic-equiv-3} there is an isomorphism of monads \(Q(H) \cong D(H)\).
  \end{thmlist}
  Additionally, if \(\cat{C}\) is pivotal with pivotal structure \(\phi\),
  any of the above statements hold if and only if \(H\) and \(\phi\) admit a pair in involution.
\end{theorem}
\begin{proof}\ref{big-monadic-equiv-1}\(\implies\)\ref{big-monadic-equiv-2}:
  Suppose \(\omega \in \QCat*{\cat{C}^H}\) with \(U^{Q(H)}\omega = 1\).
  As shown in \cref{eq: objs_as_mod_funs},
  this induces a functor of module categories
  \[
    \omega \otimes \blank  \from \cat{C}^{D(H)} \to \cat{C}^{Q(H)}.
  \]
  Since \(U^{Q(H)}\omega = 1 \in \cat{C}\),
  we can apply \cref{prop: functors_as_morphisms_of_monads},
  and obtain that \(Q(H)\) and \(D(H)\) are isomorphic as comodule monads.

  It immediately follows that~\ref{big-monadic-equiv-2} implies~\ref{big-monadic-equiv-3};
  we proceed with~\ref{big-monadic-equiv-3}\(\implies\)\ref{big-monadic-equiv-1}:
  consider an isomorphism of monads \(g \from Q(H) \nt D(H)\).
  It gives rise to a functor \(G \from \cat{C}^{D(H)} \to \cat{C}^{Q(H)}\) that, on objects, is defined by
  \[
    G(m, \nabla_m) = (m, \nabla_m g_m), \qquad \qquad \text{ for all } (m, \nabla_m) \in \cat{C}^{D(H)}.
  \]
  We compose \(G\) with \(E^{Q(H)} \from \cat{C}^{Q(H)} \to \QCat*{\cat{C}^H}\)—%
  the inverse of the comparison functor defined in Equation~\eqref{eq: inverse_comparison_functor}—%
  and see that there exists an object
  \[
    1^{(Q)} \defeq E^{Q(H)}G(1) \in \QCat*{\cat{C}^H},
  \]
  whose underlying object is the unit of \(\cat{C}\).
  Now let \((\cat{C}, \phi)\) be pivotal.
  By \cref{lem: quasi_piv_vs_inv_ayds}, lifts of \(1\in \cat{C}\) to the dual of
  the anti\hyp{}centre \(\QCat*{\cat{C}^H}\)
  are in correspondence with quasi-pivotal structures
  \((\beta, \rho_{\beta})\), where \(\beta \in \Char(H)\).
  By \cref{prop: piis_are_qpiv}, such a quasi-pivotal structure exists if and only if
  there exists a pair in involution for \(H\) and \(\phi\).
\end{proof}

As a corollary, we can determine whether a category is pivotal
in terms of monad isomorphisms between the central and anti\hyp{}central monad.
For a category \(\cat{C}\),
recall \cref{def: central_bimonad} of its central Hopf monad,
and \cref{def: anti_central_comodule_monad} of its anti\hyp{}central comodule monad.

\begin{corollary}\label{cor: pivotal_from_central_anti_central}
  Let \(\cat{C}\) be a rigid monoidal category.
  If \(\cat{C}\) admits a central Hopf monad \(\mathfrak{D}(\cat{C})\)
  and an anti\hyp{}central comodule monad \(\mathfrak{Q}(\cat{C})\),
  then it is pivotal if and only if \(\mathfrak{D}(\cat{C}) \cong \mathfrak{Q}(\cat{C})\) as monads.
\end{corollary}
\begin{proof}
  We consider the identity \(\Id_{\cat C} \from \cat C \to \cat C\) as a Hopf monad.
  Its Drinfeld and anti\hyp{}Drinfeld double are
  \(D(\Id_{\cat C}) = \mathfrak D(\cat C) \rtimes \Id_{\cat C}\)
  and \(Q(\Id_{\cat C}) = \mathfrak Q(\cat C) \rtimes \Id_{\cat C}\).
  From here it follows that \(D(\Id_{\cat C}) = \mathfrak D(\cat C)\),
  and similarly \(Q(\Id_{\cat C}) = \mathfrak Q(\cat C)\).
  The proof is concluded by \cref{thm: equivalence_in_the_monad_setting}.
\end{proof}

%%% Local Variables:
%%% mode: LaTeX
%%% TeX-master: "../main"
%%% TeX-command-extra-options: "-shell-escape -fmt=prec.fmt -file-line-error -synctex=1"
%%% End:
